\section{The odd conormal is the original strict Selmer dual}
\label{sec:kt-projector-equivalence}

This is a scope audit of the new projector criterion, not an
additional arithmetic vanishing argument.  We retain the
hypotheses and the actual normalized ring $S$ of
KT-EQUIVARIANT-MARKED-SLICE.  The source's rational trivial
character is the character of modulus $N=2$:
\texttt{SmallClassNumbers/ClassNumberThree.lean:13--32}
identifies its module with
\[
 \mathcal X=\operatorname{Sha}_{2\ell}
              (G_{\mathbf Q},(K_O/O)(\alpha)),\qquad
 K_O=\operatorname{Frac}O,\quad
 \alpha=\Psi^{-1}=\varepsilon_\ell^{-e}.                    \tag{KE.1}
\]
Here $\operatorname{Sha}_{2\ell}$ means exactly the source's
\texttt{ShaU}: the cohomology class restricts to zero on
every inertia group and on every decomposition group over
primes dividing $2\ell$.  The literal definition is
\texttt{GaloisLibrary/SelmerTransport/ConjStableShaFam.lean:94--96,167--169}.
The identification of the coefficient module
$\texttt{QO}=O\otimes_{\mathbf Z_\ell}(\mathbf Q_\ell/\mathbf Z_\ell)$
with $K_O/O$ follows by tensoring with the finite free
$\mathbf Z_\ell$-module $O$.

\begin{theorem}[KT-ODD-CONORMAL-SELMER-EQUIVALENCE]
Put
\[
 C=\bigl(\ker(x_0:S\to O)/(\ker x_0)^2\bigr)_{-1}^{\kappa_S=-1}.
                                                                  \tag{KE.2}
\]
For every $0\ne b\in O$ there is an isomorphism
\[
        \operatorname{Hom}_O(C,O/(b))\simeq\mathcal X[b].    \tag{KE.3}
\]
Use $O/(b)\simeq(K_O/O)[b]$, $u\mapsto u/b$, to specify
the coefficient identification.  The isomorphisms are
compatible with the usual inclusions of torsion chunks;
in particular
\[
             \mathcal X\simeq\operatorname{Hom}_O(C,K_O/O).
                                                                  \tag{KE.4}
\]
Thus $C=0$ if and only if $\mathcal X=0$.  Under the
established finiteness of $\mathcal X$ in the carried
source chain, $C$ has finite length and
\[
                  \operatorname{length}_O C
                       =\operatorname{length}_O\mathcal X. \tag{KE.5}
\]
Consequently the unit-minor criterion of
KT-ODD-CONORMAL-GATE is equivalent to the original strict
Selmer vanishing problem.  Without an independently
evaluated minor or a new arithmetic comparison, it does
not constitute a stricter arithmetic gate or the removal
of an automorphic branch.
\end{theorem}

\begin{proof}
Write $M=(K_O/O)(\alpha)$ and $A_b=(O/(b))(\alpha)$.
We give both directions, including the coefficient and
local-condition checks.

\paragraph{Global coefficient identification and normalization.}
The marked element $g_1$ satisfies
$\Psi(g_1)^2-1\in O^\times$.  Therefore
$\alpha(g_1)-1\in O^\times$, so $M^{G_{\mathbf Q}}=0$.
The exact sequence
\[
       0\longrightarrow A_b\longrightarrow M
         \xrightarrow{\ b\ }M\longrightarrow0              \tag{KE.6}
\]
gives
$H^1(G_{\mathbf Q},A_b)\simeq H^1(G_{\mathbf Q},M)[b]$.
Every class on the left has a unique cocycle representative
$t$ with $t(g_1)=0$: subtract its uniquely determined
coboundary using the unit $\alpha(g_1)-1$.  This is the
normalization performed in
\texttt{ShaInjGen.shaInj\_norm:33--99}.  Also $t(c)=0$,
because $c^2=1$, $\alpha(c)=1$, and two is invertible.

\paragraph{From an odd conormal functional to a scalar cocycle.}
A functional in $\operatorname{Hom}_O(C,O/(b))$ extends by
zero on every other weight and sign to an $O$-derivation
$D:S\to O/(b)$ at $x_0$.  Conversely such a derivation
factors uniquely through $C$.  The homomorphism
\[
           F(s)=x_0(s)+\epsilon D(s),\qquad
           \epsilon^2=0,\quad b\epsilon=0                  \tag{KE.7}
\]
takes values in the complete local ring
$O\oplus\epsilon O/(b)$, which is \texttt{IsCO}.
Continuity introduces no extra condition: if
$\mathfrak m_S=\mathfrak m_O S+\ker x_0$, Leibniz gives
$D(\mathfrak m_S^r)\subset\mathfrak m_O^{r-1}O/(b)$,
which is zero for all sufficiently large $r$.
The universal representation evaluated at (KE.7) is a
continuous dual-number lift.

Pass through the isomorphism to the genuinely equivariant
$F$-marked slice of KT-EQUIVARIANT-MARKED-SLICE.  On this
slice $\kappa$ is the literal matrix operation (KP.6).
The left multiplicative cocycle of the lift therefore has
weight $-1$ and is odd for
$X\mapsto-KX^tK$.  Hence it has exactly the form
\[
         \gamma(g)=t(g)Y_-,\qquad Y_-=E_{21}-E_{32},        \tag{KE.8}
\]
and the scalar cocycle relation is
$t(gh)=t(g)+\alpha(g)t(h)$, by (HS.2).
At $g_1$, put $q=\Psi(g_1)$.  The linearized second odd
marking (KP.9), applied to (KE.8), is
\[
                         (q+1)t(g_1)=0.                    \tag{KE.9}
\]
Since $q+1$ is a unit, $t(g_1)=0$.  The old marking has
the same linearized value on this odd direction, so both
Cayley changes of slice are the identity for this
dual-number lift.  Thus (KE.8) and its normalization
describe the functional on the actual original ring as
well; we have not confused a transported tangent sign
with a literal matrix-entry equation.

\paragraph{The local condition at $\ell$ is strict.}
KT-ORDINARY-INTEGRAL-COMPONENT identifies the actual local
quotient $U.R/I_t$ with the integral ordered ordinary
functor.  Therefore the local lift just constructed has
its unique ordinary flag.  After setting $\epsilon=0$
this is the standard flag of $r_0$: its ordered inertial
eigenvalues are distinct.  Its lower-unipotent flag
coordinate consequently has the form $1+\epsilon W$,
with $W$ strictly lower triangular over $O/(b)$.
Writing the lift as the conjugate of an upper Borel
lift and comparing lower entries gives
\[
       \gamma|_{G_{\mathbf Q_\ell}}
            =\operatorname{cob}(W)
                       \quad\text{in the lower-weight space}.
                                                                  \tag{KE.10}
\]
The base adjoint action preserves weights and commutes
with $\kappa_{\rm Lie}$.  Project (KE.10) to weight $-1$
and sign $-1$.  The projected $W$ is a scalar multiple
of $Y_-$, so $t|_{G_{\mathbf Q_\ell}}$ is a coboundary
already with coefficients in $A_b$.  This argument works
for both split and unramified nonsplit $\eta$.

\paragraph{All away-prime inertia conditions.}
For the trivial Dirichlet character, the exact
\texttt{DatFacts.dat\_facts} description of
\texttt{Kill} is the union, over all places away from
$\ell$, of
\[
                        I_q\cap\ker\eta .                  \tag{KE.11}
\]
See \texttt{RhoKApply.lean:61--62} and
\texttt{DatFacts.lean:70--100}.
The cyclotomic character is unramified at each such
place, so $\alpha|_{I_q}=1$ and $t|_{I_q}$ is a
continuous additive homomorphism.  The box condition
makes it zero on (KE.11).  It therefore factors through
a group of order at most two.  Two is a unit on $A_b$,
so this homomorphism is zero on all of $I_q$.  This
includes the ramified quadratic place $q=3$; no
unramifiedness of $\eta$ at that place is being assumed.
Together with the local strict condition at $\ell$,
all inertia groups in the literal \texttt{shaFam} have
been handled.

\paragraph{The extra decomposition group at $2$.}
The parameter is $2\ell$, so the source also requires
triviality on a full decomposition group at $2$.
After embedding $A_b$ in $M$, the class just obtained
is unramified there.  Its unramified cohomology is
\[
 H^1_{\rm unr}(G_{\mathbf Q_2},M)
       =M/(2^{-e}-1)M=0.                                   \tag{KE.12}
\]
Indeed $e>0$, so $2^{-e}-1\ne0$ in $K_O$, and
multiplication by any nonzero element of $O$ is
surjective on $K_O/O$.  Formula (KE.12) uses divisible
coefficients.  The corresponding group with coefficients
in $O/(b)$ need not vanish, and no such assertion is
used.  The usual conjugacy of decomposition groups
handles every extension of the place to the algebraic
closure.  The definitions of \texttt{shaFam} and
\texttt{ShaU} now show that the image lies in
$\mathcal X[b]$.

\paragraph{The inverse construction.}
Given $\xi\in\mathcal X[b]$, use (KE.6) to obtain its
unique finite-coefficient class and normalize its
cocycle $t$ at $g_1$.  On every inertia group away from
$\ell$ the action on $A_b$ and $M$ is trivial.  Since
$\xi$ restricts to zero in $H^1(I_q,M)$ and the
coefficient inclusion is injective, $t|_{I_q}=0$.

At $\ell$, choose an inertia element $\tau$ for which
$\bar\Psi(\tau)\ne1$, as supplied by GD.  Then
$\alpha(\tau)-1$ is a unit and
$M^{G_{\mathbf Q_\ell}}=0$.  Applying (KE.6) locally
shows that
\[
 H^1(G_{\mathbf Q_\ell},A_b)
                   \longrightarrow H^1(G_{\mathbf Q_\ell},M)
                    \quad\text{is injective}.              \tag{KE.13}
\]
The strict local condition on $\xi$ therefore makes
$t$ a local coboundary with finite coefficients.  Form
the cocycle (KE.8) and the corresponding dual-number
GG lift.  It kills \texttt{Kill}.  Its local restriction
is conjugate by $1+\epsilon wY_-$ to the exact base
representation, hence is integrally ordinary with the
specified inertia; the local component lemma gives
the required $I_t$ condition.  The equations
$t(c)=t(g_1)=0$ give the source normalization.  The
universal property of $S$ now supplies (KE.7).
It is odd and of weight $-1$, so its derivation factors
through $C$.

Uniqueness of the $g_1$-normalized scalar cocycle and
uniqueness in \texttt{S.rep} show that these two
constructions are inverses.  They respect addition
and scalar multiplication.  The choice of cyclic
generator in the archived \texttt{sha\_inj} may change
its displayed cocycle map by a coefficient unit;
the explicit identification $u\mapsto u/b$ fixes that
choice in (KE.3).

Finally use $b=\varpi^n$ and take the direct limit.
Every homomorphism from the finite $O$-module $C$ to
$K_O/O$ has bounded torsion image, giving (KE.4).
This dual functor detects every nonzero finite
$O$-module.  Its elementary-divisor decomposition
shows that it has a divisible summand for every free
summand of $C$, while a summand $O/(\varpi^a)$ is
dual to another summand of the same length.  Finiteness
of $\mathcal X$ therefore forces $C$ to be torsion
and proves (KE.5).
\end{proof}
